\chapter{State Diagram}
\label{ch:state}

\section{Purpose}

The state diagram models the lifecycle of a tool session: the states the tool
occupies, the events that move it between them, and the point at which the user
decides whether anything is written.

It models the tool's own session, not the behaviour the tool translates. The
behavioural models the tool produces are the subject of
Chapter~\ref{ch:activity}.

\umlfigure{state-diagram}{State diagram for a tool session}{fig:state}

\section{States}

\begin{longtable}{L{0.26\textwidth}L{0.66\textwidth}}
	\caption{Session states}\label{tab:states}\\
	\toprule
	\textbf{State} & \textbf{Description} \\
	\midrule
	\endfirsthead
	\toprule
	\textbf{State} & \textbf{Description} \\
	\midrule
	\endhead

	Idle & The initial state. The tool is not running and GENESYS is in normal
	use. \\
	\addlinespace
	Waiting for User Configuration & The dialog is open and the engineer is
	choosing a translation direction and its parameters. The tool waits
	indefinitely. \\
	\addlinespace
	Collecting Input & The tool gathers the material the translation needs,
	reading the selected source files or the selected model. \\
	\addlinespace
	Analyzing & Extraction, abstraction and mapping run. This is the only state
	with two outcomes. \\
	\addlinespace
	Awaiting Review & The proposed change set is displayed and the tool waits for
	the engineer. Nothing has been written at this point. \\
	\addlinespace
	Writing & The accepted changes are applied, either to the GENESYS model or to
	the code repository. \\
	\addlinespace
	Showing Report & Results or errors are displayed. Reached from a successful
	write, from a denied review, and from a failed analysis. \\
	\bottomrule
\end{longtable}

\section{Transitions}

\begin{longtable}{L{0.26\textwidth}L{0.26\textwidth}L{0.38\textwidth}}
	\caption{Session transitions}\label{tab:transitions}\\
	\toprule
	\textbf{From} & \textbf{To} & \textbf{Trigger} \\
	\midrule
	\endfirsthead
	\toprule
	\textbf{From} & \textbf{To} & \textbf{Trigger} \\
	\midrule
	\endhead
	Initial & Idle & --- \\
	Idle & Waiting for User Configuration & Click on Toolbar \\
	Waiting for User Configuration & Collecting Input & Run \\
	Collecting Input & Analyzing & Done \\
	Analyzing & Awaiting Review & Success \\
	Analyzing & Showing Report & Error \\
	Awaiting Review & Writing & User Accepts \\
	Awaiting Review & Showing Report & User Denies \\
	Writing & Showing Report & Done \\
	Showing Report & Idle & Close \\
	\bottomrule
\end{longtable}

\section{Design Decisions}

\paragraph{Two states are actually waits.}
Waiting for User Configuration and Awaiting Review are states rather than steps,
because the tool stops indefinitely in each until the engineer acts. Every other
state could be interpretted as an action and is transient and leaves of its own accord.

\paragraph{Nothing is written before review.}
Writing is reachable only through Awaiting Review and only on acceptance. Every
other path to Showing Report leaves the model and the repository untouched, so
the guarantee that a cancelled or failed session writes nothing is a property of
the diagram's shape rather than a note attached to it.

\paragraph{One reporting state for three outcomes.}
Success, denial and failure all converge on Showing Report. This keeps the
diagram compact, and means the engineer always returns to Idle by the same
route.

\paragraph{The decision node after review.}
Accept and deny are drawn through an explicit decision node rather than as two
transitions leaving the state directly, which makes the branch visible at a
glance.
